%% Copyright 2026
%% Licensed under Creative Commons Attribution–ShareAlike 4.0 International (CC BY-SA 4)

\documentclass[12pt]{article}

% \usepackage[cm]{fullpage}
\usepackage{amssymb}
\usepackage{amsmath}
\usepackage{tcolorbox}
\usepackage[makeroom]{cancel}
\usepackage{tikz}

\title{Solution Notes for 12-Ball Problem}
\date{June 1, 1999}

% no indent for whole document
\setlength\parindent{0pt}

\begin{document}
\maketitle

\section*{The 12-ball  Problem}

\begin{itemize}
\item 12 balls are identical in every way, except one is either lighter or heavier than the rest  (Think featureless billiard balls).  Among the 12 balls is an {\textbf{odd ball}}.
\item One of the balls is odd -- it's either {\textbf{lighter}} or
{\textbf{heavier}} than the other 11 balls. You cannot tell which ball is odd
by visual inspection.  
\item You must use a simple balance scale to identify the {\textit{odd ball}} AND
determine what {\textit{kind}} of ball it is (light or heavy).

\end{itemize}

\subsection*{Limits}
\begin{enumerate}
\item You're allowed to use a balance-scale to weigh any the balls in any
 combination you wish.  For sake of clarity, each application of the scale
is called a {\textbf{weighing}}.  You are allowed only {\textbf{three}} weighings.
\item You are allowed to place any balls from the set of 12 balls on the balance scale only.
\end{enumerate}

\section*{Soluition}

The following pages describe the solution.  {\textsc{Stop reading now}} unless
you want the solution given to you.


\newpage
\section*{Solution}


First some notation:\\[2mm]

\begin{itemize}
\item Let the balls be a set $B = { b_1, b_2, b_3, \ldots, b_{12}}$
\item This is the notation that ball $b_1$ is ball 1, and $b_{9}$ is ball 9, and so on.
\item Each weighing event is identified as $W_1$, $W_2$, and $W_3$.
\item If we know a ball is normal, we just refer to it as $b_N$. It's identity is meaningless.
\item If we know a ball is heavy (and not light), we refer to it as $b_{\textrm{heavy}}$.
\item If we know a ball is light (and not heavy), we refer to it as $b_{\textrm{light}}$.

\end{itemize}

\subsection*{First Weighing $W_1$}

$W_1$ = (left-side) ${ b_1, b_2, b_3, b_4}$ against (right-side) ${b_5, b_6, b_7, b_8}$.

If the scale is even, then record Possibly Heavy (PH), Possibly Light (PL), and Normal (N) as:

\begin{align*}
N &= {b_1, b_2, b_3, b_4, b_5, b_6, b_7, b_8}\\
PH  &= {b_9, b_{10}, b_{11}, b_{12}}\\
PL  &= {b_9, b_{10}, b_{11}, b_{12}}
\end{align*}

If the scale is heavy left,  then record:

\begin{align*}
PH &= {b_1, b_2, b_3, b_4}\\
PL &= {b_5, b_6, b_7, b_8}\\
N  &= {b_9, b_{10}, b_{11}, b_{12}}
\end{align*}

If the scale is heavy right, then record:

\begin{align*}
PH &= {b_5, b_6, b_7, b_8}\\
PL &= {b_1, b_2, b_3, b_4}\\
N  &= {b_9, b_{10}, b_{11}, b_{12}}
\end{align*}

\subsection*{Second Weighing $W_2$}

\subsubsection*{$W_1$ was balanced}

If $W_1$ was balanced, then the odd ball is among ${b_9, b_{10}, b_{11}, b_{12}}$, then\\

$W_2$ = (left-side) ${b_N, b_N, b_9, b_{10}}$ against (right-side) ${b_N, b_N, b_{11}, b_{12}}$.\\

Observe which way the scale moves.\\

If it is heavy-left, then record:
\begin{align*}
PH &= {b_9, b_{10}}\\
PL &= {b_{11}, b_{12}}
\end{align*}

If it is heavy-right, then record:
\begin{align*}
PL &= {b_9, b_{10}}\\
PH &= {b_{11}, b_{12}}
\end{align*}

It cannot be balanced since the balls ${b_9, \ldots, b_{12}}$ did not participate in the even weighing event $W_1$.  Go to ${W_3}$.


\subsubsection*{$W_1$ was heavy left}

From $W_1$ we know this:

\begin{align*}
PH &= {b_1, b_2, b_3, b_4}\\
PL &= {b_5, b_6, b_7, b_8}\\
N  &= {b_9, b_{10}, b_{11}, b_{12}}
\end{align*}

We will isolate one ball from each side.\\

$W_2$ = (left-side) ${b_9, b_{10}, b_4, b_5}$ against (right-side) ${b_{11}, b_{12}, b_3, b_6}$.\\

If heavy left, then record:

\begin{align*}
PH &= {b_4}\\
PL &= {b_6}\\
N  &= {b_1, b_2, b_3, b_5, b_7, b_8, b_9, b_{10}, b_{11}, b_{12}}
\end{align*}

If heavy right, then record:
\begin{align*}
PH &= {b_3}\\
PL &= {b_5}\\
N  &= {b_1, b_2, b_4, b_6, b_7, b_8, b_9, b_{10}, b_{11}, b_{12}}
\end{align*}

\subsubsection*{$W_1$ was heavy right}

From $W_1$ we know this:

\begin{align*}
PH &= {b_5, b_6, b_7, b_8}\\
PL &= {b_1, b_2, b_3, b_4}\\
N  &= {b_9, b_{10}, b_{11}, b_{12}}
\end{align*}

We will isolate one ball from each side. \\
$W_2$ = (left-side) ${b_9, b_{10}, b_8, b_1}$ against (right-side) ${b_{11}, b_{12}, b_7, b_2}$.\\


If heavy-left, we record:
\begin{align*}
PH &= {b_8}\\
PL &= {b_5}\\
N  &= {b_1, b_2, b_4, b_6, b_7, b_8, b_9, b_{10}, b_{11}, b_{12}}
\end{align*}



\subsection*{Third Weighing $W_3$}

\subsubsection*{$W_2$ was heavy-left}

If $W_2$ was heavy-left, then:\\

${W_3}$ = (left-side) ${b_N, b_N, b_9, b_{11}}$ against (right-side) ${b_N, b_N, b_N, b_N}$.\\

Observe the scale.\\

If $W_3$ is heavy-left, then $b_{\textrm{heavy}} = b_9$. ($b_{11}$ cannot be both heavy and light)\\
If $W_3$ is heavy-right, then $b_{\textrm{light}} = b_{11}$ ($b_9$ cannot be both heavy and light,
and $b_{11}$ moved right to left)\\

\subsubsection*{$W_2$ was heavy-right}

If $W_2$ was heavy-right, then:\\

$W_3$ = (left-side) ${b_N, b_N, b_{10}, b_{12}}$ against (right-side) ${b_N, b_N, b_N, b_N}$.\\
If $W_3$ is heavy-left, then $b_{\textrm{heavy}} = b_{10}$. ($b_{12}$ cannot be both heavy and light)\\
If $W_3$ is heavy-right, then $b_{\textrm{light}} = b_{12}$ ($b_{12}$ moved from right to left)\\



\end{document}

